We organise the six metrics into a decision table by use case.
The table answers the practitioner question: \emph{given a specific legal or
analytical task, which metric should I lead with?}

\begin{table}[ht]
\centering
\caption{Compactness metric decision table for redistricting practitioners.}
\label{tab:decision}
\begin{tabular}{p{3.0cm}p{2.2cm}p{5.5cm}p{3.0cm}}
\toprule
Use Case & Lead Metric & Rationale & \textsc{bisect} Flag \\
\midrule
Litigation evidence (federal/state court) &
  PP &
  Most frequently cited metric in U.S.\ redistricting litigation; $[0,1]$ scale
  explainable to lay juries in minutes; PP $= 4\pi A/P^2$ derivable on a whiteboard &
  \texttt{--types pp} \\
\addlinespace
Coastal/riverine state comparison &
  Reock &
  Less sensitive to TIGER/Line boundary artifacts than PP; enclosing-circle visual
  intuitive; disclose \textsc{bisect}'s centroid-radius proxy when exact MBC matters &
  \texttt{--types reock} \\
\addlinespace
Public communication &
  CH &
  Most visually intuitive: audience can see whether district fits inside its convex
  hull without mathematical training; effective for showing non-convex appendages &
  \texttt{--types ch} \\
\addlinespace
VRA district analysis &
  PWC &
  Measures representational compactness; compact geometry does not guarantee
  that minority voters are efficiently clustered around the district centroid &
  \texttt{--types pwc} \\
\addlinespace
Statutory compliance (CO, IA, OR, UT) &
  S &
  Colorado's constitution references a Schwartzberg-equivalent criterion;
  $S = 1/\sqrt{\text{PP}}$ allows conversion from PP output to S for filings &
  \texttt{--types schwartzberg} \\
\addlinespace
Algorithmic optimisation target &
  LW &
  LW $>3$ threshold is the most commonly enforced bright-line rule in court
  filings; bisect produces LW $\leq 2.8$ on NC 2020 by construction &
  \texttt{--types lw} \\
\bottomrule
\end{tabular}
\end{table}

\subsection{Notes on the Decision Table}

\paragraph{Lead metric is not the only metric.}
The decision table identifies the \emph{primary} metric for each use case.
We recommend always including all six metrics in a certified compactness report
(K.7), using the lead metric to organise testimony or narrative while
supplementing with the full profile.

\paragraph{Cross-family redundancy.}
Because PP and S are algebraically equivalent ($S = 1/\sqrt{\text{PP}}$),
a report that includes PP need not separately report S unless the jurisdiction's
statute explicitly requires S notation. The conversion table in K.4 enables
translation between representations.

\paragraph{Algorithm selection.}
The decision table does not depend on algorithm choice for the lead metric
recommendation. However, if the goal is \emph{maximising} a specific metric,
algorithm matters:
\begin{itemize}
  \item Improve the Reock-oriented proxy: use \texttt{--structure moving-knife}
    with the K.2 disclosure that this is not exact polygon-MBC optimisation.
  \item Minimise PWC: use \texttt{--structure prime-factor} (K.6).
  \item Maximise PP: use \texttt{--structure ratio-optimal} (K.1).
  \item All algorithms guarantee CH $> 0.85$ on NC by construction (K.3).
\end{itemize}
